\documentclass[leqno]{article}
% Shared mark scheme preamble for resources.danieljsmith.org: the pack
% preamble, unchanged, then the scheme's own layer. A scheme's .tex is
% rendered from its YAML (scripts/README.md, Mark schemes), never edited.
\input{\CurrentFilePath/pack_preamble.tex}

% The page: equal side margins, so the table runs the full width of the
% header rule. The footers' right edge follows the text's.
\geometry{left=0.8in, right=0.8in}
\makeatletter
\renewcommand{\djs@footR}[1]{%
  \put(\dimexpr1in+\hoffset+\oddsidemargin+\textwidth\relax,
    -\dimexpr\paperheight-0.32in\relax){\makebox[0pt][r]{#1}}}
\makeatother

\usepackage[export]{adjustbox}
\usepackage{tabularray}
\definecolor{msGrid}{RGB}{176, 172, 164}
\definecolor{msTint}{RGB}{245, 243, 239}

% The header on every page, the front page included, and the booklet label
% under the contents-page title. \djsQuestionPageStyle stays empty.
\newcommand{\djsSchemeHeader}[1]{%
  \djsPackHeader[markscheme]{\textit{Mark scheme} - #1}%
  \renewcommand{\djsPackTitle}{#1}%
  \renewcommand{\djsBookletLabel}{Mark scheme}%
}

% The front page: the contents heading, spaced as tocloft spaces the packs'
% "Questions", then the contents table. A question's name and page number
% both link to its first page.
\newcommand{\msContentsHeading}{%
  \vskip\cftbeforetoctitleskip
  {\noindent\cfttoctitlefont\contentsname\par}%
  \vskip\cftaftertoctitleskip}
\NewTblrEnviron{mscontents}
\SetTblrInner[mscontents]{
  colspec = {Q[wd=26mm, l] X[l] Q[wd=16mm, c] Q[wd=14mm, r]},
  width = \linewidth,
  rowsep = 4pt, colsep = 6pt,
  hlines = {0.4pt, msGrid}, vlines = {0.4pt, msGrid},
  hline{2} = {2.2pt, scarlet},
  row{1} = {font=\small\bfseries},
  row{Z} = {font=\bfseries},
  column{2} = {font=\small},
}
\newcommand{\msContentsQuestion}[1]{\hyperref[q:#1]{Question #1}}
\newcommand{\msContentsPage}[1]{\hyperref[q:#1]{\pageref*{q:#1}}}
\newcommand{\msSep}{\hspace{0.9em}}

% The conventions pages' layout: headings in the documents' own bold, tables
% ruled and padded as the scheme's, roomy lists.
\newcommand{\msHead}[1]{%
  \par\vspace{18pt}%
  {\noindent\bfseries #1\par}%
  \vspace{6pt}}
\NewTblrEnviron{mstblr}
\SetTblrInner[mstblr]{rowsep = 6pt, colsep = 6pt,
  hlines = {0.4pt, msGrid}, vlines = {0.4pt, msGrid}}
\newenvironment{msList}
  {\begin{itemize}[leftmargin=1.2em, itemsep=5pt, topsep=3pt, parsep=0pt]}
  {\end{itemize}}

% A question's first page: the link target and the PDF bookmark, as
% \questionitem gives the packs.
\newcommand{\msQnum}{}
\newcommand{\msQuestion}[1]{%
  \clearpage
  \renewcommand{\msQnum}{#1}%
  \phantomsection\label{q:#1}%
  \addcontentsline{toc}{subsection}{\protect\djsTocEntry{#1}}}

% The scheme's table: Question 16mm, Scheme the rest, Marks 16mm, Notes
% 44mm; every cell ruled in msGrid; the scarlet rule under the head row. A
% long table has no caption or "continued" foot; a continuation page opens
% with a muted line above the repeated head row. The rules across a part's
% rows, the alternative's tint and the total row are the rendered table's.
\NewTblrTheme{ms}{%
  \DefTblrTemplate{firsthead}{default}{}%
  \DefTblrTemplate{middlehead,lasthead}{default}{%
    {\small\itshape\color{djsMuted}Question \msQnum\ continued\par}%
    \vspace{6pt}}%
  \DefTblrTemplate{firstfoot,middlefoot,lastfoot}{default}{}%
}
\NewTblrEnviron{msscheme}
\SetTblrOuter[msscheme]{long, theme=ms, entry=none, label=none}
\SetTblrInner[msscheme]{
  colspec = {Q[wd=16mm, l] X[l] Q[wd=16mm, c] Q[wd=44mm, l]},
  width = \linewidth,
  rowsep = 6pt, colsep = 6pt,
  rowhead = 1,
  vlines = {0.4pt, msGrid},
  hline{1} = {0.4pt, msGrid}, hline{2} = {2.2pt, scarlet},
  hline{Z} = {0.4pt, msGrid},
  column{1} = {font=\bfseries}, column{3} = {font=\bfseries},
  column{4} = {font=\small},
  row{1} = {rowsep=4pt, font=\small\bfseries},
}

% Two marks on a row: each note starts with its code.
\newcommand{\msNote}[2]{\textbf{#1}\enspace #2}
% A full-row figure: the worked solution's own picture, scaled uniformly and
% only down, to the table less the Question column. The picture is drawn
% with the line width it has in the solutions document, at its list depth
% there (the optional argument, 1 by default: a part's working), so one
% sized by \linewidth keeps its size; tabularray sets the cell's \linewidth
% to the cell's measured width, which would shrink it at every pass. The
% widths are measured once, in a solution box at the packs' text width
% (pack_preamble.tex's margins, 0.8in and 1in).
\newlength{\msSpanW}
\makeatletter
\newcommand{\ms@keepLW}[1]{%
  \expandafter\xdef\csname ms@solnLW@#1\endcsname{\the\linewidth}}
\AtBeginDocument{%
  \setlength{\msSpanW}{\dimexpr\textwidth-16mm-28pt\relax}%
  \setbox\z@\vbox{%
    \begin{minipage}{\dimexpr\paperwidth-1.8in\relax}
      \begin{djsSolution}
        \ms@keepLW{0}%
        \begin{enumerate}\item \ms@keepLW{1}%
        \begin{enumerate}\item \ms@keepLW{2}%
        \begin{enumerate}\item \ms@keepLW{3}%
        \end{enumerate}\end{enumerate}\end{enumerate}
      \end{djsSolution}
    \end{minipage}}}
\newcommand{\msFigure}[2][1]{%
  \adjustbox{max width=\msSpanW, max totalheight=0.6\textheight, valign=t}{%
    \hsize=\csname ms@solnLW@#1\endcsname
    \linewidth=\hsize
    #2}}
\makeatother
% A figure mark's feature, in the Scheme column.
\newcommand{\msFeature}[1]{\textit{#1}}
% An alternative's title row.
\newcommand{\msAlternative}[1]{\textit{Alternative:} #1}
% Drafts only: a queried mark's code, the query's why under its note, and a
% flag in its part's place.
\newcommand{\msQueried}[2]{#1~\(\to\) #2?}
\newcommand{\msQuery}[1]{{\itshape\color{djsMuted}Query: #1}}
\newcommand{\msFlag}[2]{\textbf{#1:}\enspace #2}


\djsSchemeHeader{Power and Driving Force}
\djsSitePath{Further Mechanics}{Work, Energy and Power}

\djsLastUpdated{30 October 2026}
\begin{document}
\djsSetQuestionCount{12}
\djsAddMarks{1}{8}
\djsAddMarks{2}{5}
\djsAddMarks{3}{14}
\djsAddMarks{4}{11}
\djsAddMarks{5}{7}
\djsAddMarks{6}{10}
\djsAddMarks{7}{7}
\djsAddMarks{8}{8}
\djsAddMarks{9}{9}
\djsAddMarks{10}{9}
\djsAddMarks{11}{8}
\djsAddMarks{12}{11}
\vspace*{6pt}
\djsContentsTitle
\vspace{30pt}
\msContentsHeading
\begin{mscontents}{}
Question & Marks by part & Marks & Page \\
\msContentsQuestion{1} & (a)\,2\msSep (b)\,2\msSep (c)\,4 & 8 & \msContentsPage{1} \\
\msContentsQuestion{2} & & 5 & \msContentsPage{2} \\
\msContentsQuestion{3} & (i)\,2\msSep (ii)\,2\msSep (iii)(a)\,4\msSep (iii)(b)\,2\msSep (iv)(a)\,2\msSep (iv)(b)\,2 & 14 & \msContentsPage{3} \\
\msContentsQuestion{4} & (a)\,4\msSep (b)\,7 & 11 & \msContentsPage{4} \\
\msContentsQuestion{5} & & 7 & \msContentsPage{5} \\
\msContentsQuestion{6} & (a)\,4\msSep (b)\,2\msSep (c)\,2\msSep (d)\,2 & 10 & \msContentsPage{6} \\
\msContentsQuestion{7} & & 7 & \msContentsPage{7} \\
\msContentsQuestion{8} & & 8 & \msContentsPage{8} \\
\msContentsQuestion{9} & (a)\,4\msSep (b)\,5 & 9 & \msContentsPage{9} \\
\msContentsQuestion{10} & (a)\,4\msSep (b)\,5 & 9 & \msContentsPage{10} \\
\msContentsQuestion{11} & & 8 & \msContentsPage{11} \\
\msContentsQuestion{12} & (a)\,2\msSep (b)\,4\msSep (c)\,5 & 11 & \msContentsPage{12} \\
\SetCell[c=2]{r} Total & & 107 & \\
\end{mscontents}
\newpage
% The conventions pages, shared by every pack's mark scheme: one source of
% truth, \input after the front page, with no strand or pack logic. The keys
% on the first side, the rules on the second. The scheme layer defines
% \msHead, msList and mstblr.
\vspace*{6pt}%
{\noindent\cfttoctitlefont Reading the mark scheme\par}%
\vspace{10pt}%
These rules apply throughout, unless a note says otherwise.

\msHead{Marks}
Every mark is worth one, so \textbf{M1 A1} on a row is two marks. Only marks
coded ft follow through.\par
\vspace{10pt}%
\begin{mstblr}{
  colspec = {Q[wd=15mm, l] X[l]},
  width = \linewidth,
  hline{2} = {2.2pt, scarlet},
  column{1} = {font=\bfseries},
  row{1} = {font=\small\bfseries, rowsep=4pt},
}
Code & Meaning \\
M1 & Method: a correct method, applied to this question. A slip doesn't cost it if the method is clear. A formula must be used, not just quoted. \\
dM1 & Dependent method: as M1, but only if the method mark before it was earned, unless its note names another. \\
A1 & Accuracy: a correct answer or step. It needs the method mark it follows. \\
A1ft & As A1, following through from an earlier answer that the note names. \\
A1~AG & The last mark of a part whose answer is printed in the question: complete working, with no errors. \\
B1 & A correct value, fact or statement. It needs no method mark. \\
B1ft & As B1, following through from an earlier answer that the note names. \\
E1 & A written explanation, interpretation or conclusion. Concluding a proof by induction needs every earlier mark. \\
\end{mstblr}

\msHead{Abbreviations}
ft and AG qualify a code in the Marks column. The others appear in the
notes.\par
\vspace{10pt}%
\begin{mstblr}{
  colspec = {Q[wd=17mm, l] X[l] Q[wd=15mm, l] X[l]},
  width = \linewidth,
  column{1} = {font=\bfseries}, column{3} = {font=\bfseries},
}
ft & follow through & AG & answer given in the question \\
cao & correct answer only, with no follow-through & cso & correct solution only: no errors in the part \\
oe & or an equivalent form & awrt & any answer that rounds to this value \\
soi & seen, or implied by later working & isw & ignore working after a correct answer \\
SC & special case: credit outside the main method & their & the student's own earlier value \\
condone & a slip that doesn't cost the mark & [\;\,] & working the mark doesn't need \\
s.f. & significant figures & d.p. & decimal places \\
\end{mstblr}

\newpage
\msHead{Reading the scheme}
\begin{msList}
\item Each row is a line of working, with the marks it earns beside it.
  Where a row carries two marks, each note starts with its code.
\item \textbf{[\textit{n}]} ends each part with its total.
\item A shaded \textit{Alternative} is a second method for the same total.
  Mark an answer by the method it follows.
\item Any other valid method earns the equivalent marks, unless the question
  names a method or says \textit{hence}.
\item A sketch earns one B1 for each feature named below its figure.
\end{msList}

\msHead{Accuracy and answers}
\begin{msList}
\item Unless a question says otherwise, give answers to 3 significant
  figures, and angles in degrees to 1 decimal place. Angles in radians take
  3 significant figures.
\item A final answer with more figures is accepted if it rounds to the
  scheme's value; one with fewer is not. An answer spoiled by rounding too
  early loses its accuracy mark.
\item Where a question asks for an exact answer, a decimal loses the final
  accuracy mark. Otherwise either is accepted.
\item Wrong or missing units cost nothing unless the question asks for them.
\item A correct answer with no working earns full marks, unless the question
  needs the working: show that, prove, hence, a method it names, or an exact
  answer. A wrong answer with no working earns nothing. Calculator answers
  are treated the same way.
\item A correct answer, once reached, isn't spoiled by later working (isw).
\end{msList}

\msHead{Value of $\boldsymbol{g}$}
\begin{msList}
\item Take $g = 9.8\,\mathrm{m\,s^{-2}}$ unless a question states a value,
  which is then the only one accepted. Answers that use $g = 9.8$ may be
  given to 2 or 3 significant figures.
\item Where a question doesn't state $g$, an answer worked with $g = 10$
  throughout earns the same marks, keeping any value printed in the
  question. The worked solutions list the final answers with $g = 10$.
\end{msList}


\msQuestion{1}
\begin{msscheme}{
  hline{3-4,6-7,9-11} = {2-Z}{0.4pt, msGrid},
  hline{5,8,12} = {0.4pt, msGrid},
  row{12} = {font=\bfseries},
}
Question & Scheme & Marks & Notes \\
1(a) & {\(P = Fv\) with \(F = R\):\\[4pt] \(200\,000 = 20R \Rightarrow R = \ldots\)} & M1 & {Uses \(P = Fv\) with driving force equal to resistance at \(v = 20\). Condone power left in kW.} \\
& {\(R = 10\,000\,\mathrm{N}\)} & A1 & \\
\nopagebreak & & [2] & \\
1(b) & {\(\frac{60\,000}{8} - \their{10\,000} = 10\,000a \Rightarrow a = \ldots\)} & M1 & {Forms an equation of motion with driving force \(\frac{P}{v}\), their resistance and mass \(10\,000\). Condone sign slips and power in kW.} \\
& {\(a = -0.25\,\mathrm{m\,s^{-2}}\)} & A1 & {Or a deceleration of \(0.25\). A0 for \(a = 0.25\) with no direction.} \\
\nopagebreak & & [2] & \\
1(c) & {\(F = \frac{200\,000}{10} = 20\,000\)} & B1 & \\
& {\(\their{20\,000} = \their{10\,000} + 10\,000 \times 9.8\sin\theta\)} & M1 A1 & {\msNote{M1}{Resolves along the slope with zero acceleration, using their driving force, their resistance and a weight component. Condone sign slips and \(\sin\)/\(\cos\) confusion.}\\ \msNote{A1}{Correct equation, oe.}} \\
& {\(\sin\theta = \frac{5}{49} \Rightarrow \theta = 5.9^\circ\)} & A1 & {awrt \(5.9^\circ\), or awrt \(0.102\) radians.} \\
\nopagebreak & & [4] & \\
\nopagebreak \SetCell[c=2]{r} Total for Question 1 & & 8 & \\
\end{msscheme}

\msQuestion{2}
\begin{msscheme}{
  hline{3-6} = {2-Z}{0.4pt, msGrid},
  hline{7} = {0.4pt, msGrid},
  row{7} = {font=\bfseries},
}
Question & Scheme & Marks & Notes \\
2 & {\(3000g\sin\theta = 3000 \times 9.8 \times \frac{1}{49} = 600\)} & B1 & \\
& {\(P = Fv\)\\[4pt] \(F = \their{600} + 300 + 60v \Rightarrow P = \ldots\)} & M1 & {Uses \(P = Fv\) with driving force equal to resistance plus a weight component, at speed \(v\) or 15. Condone sign slips.} \\
& {\(P = 15(600 + 300 + 60 \times 15)\)} & A1 & {Correct expression for the power at \(v = 15\), oe.} \\
& {Power needed, \(900v + 60v^2\), increases with \(v\), so\\[4pt] the greatest constant speed is at maximum power} & E1 & {States that the engine works at maximum power at the greatest constant speed, with a reason such as the power needed increasing with \(v\).} \\
& {\(P_{\max} = 27\,000\,\mathrm{W}\)} & A1 & \\
\nopagebreak \SetCell[c=2]{r} Total for Question 2 & & 5 & \\
\end{msscheme}

\msQuestion{3}
\begin{msscheme}{
  hline{3-4,6-7,9-11,13-14,16-17,19-20} = {2-Z}{0.4pt, msGrid},
  hline{5,8,12,15,18,21} = {0.4pt, msGrid},
  row{21} = {font=\bfseries},
}
Question & Scheme & Marks & Notes \\
3(i) & {\(P = Fv\) with \(F = R\): \(72\,000 = 18R \Rightarrow R = \ldots\)} & M1 & {Uses \(P = Fv\) at \(v = 18\) with the driving force equal to the resistance.} \\
& {\(R = 4000\,\mathrm{N}\)} & A1 & \\
\nopagebreak & & [2] & \\
3(ii) & {\(\frac{72\,000}{12} - \their{4000} = \ldots\ [= 6000 - 4000]\)} & M1 & {Finds the total resistance from \(P = Fv\) at \(v = 12\) and subtracts their tug resistance.} \\
& {\(2000\,\mathrm{N}\)} & A1 & \\
\nopagebreak & & [2] & \\
3(iii)(a) & {\(F - \their{4000} - \their{2000} = 4000 \times 0.25\)\\[4pt] \(\Rightarrow F = \ldots\ [= 7000]\)} & M1 A1 & {\msNote{M1}{Forms Newton's second law for the whole system, with driving force, both their resistances and the total mass. Condone sign slips.}\\ \msNote{A1}{Correct equation, oe.}} \\
& {\(P = Fv = \their{7000} \times 10\)} & M1 & {Uses \(P = Fv\) with their driving force from Newton's second law and \(v = 10\).} \\
& {\(P = 70\,000\,\mathrm{W}\) \((70\,\mathrm{kW})\)} & A1 & \\
\nopagebreak & & [4] & \\
3(iii)(b) & {\(T - \their{2000} = 2200 \times 0.25 \Rightarrow T = \ldots\)} & M1 & {Forms Newton's second law for the aircraft alone, or the tug alone with their driving force, including the tow bar force. Condone sign slips.} \\
& {\(T = 2550\,\mathrm{N}\)} & A1 & \\
\nopagebreak & & [2] & \\
3(iv)(a) & {Lower: no resistance below \(6\,\mathrm{m\,s^{-1}}\), so greater\\[4pt] resultant force and acceleration there} & E1 & {States lower.} \\
& {Models agree from 6 to \(10\,\mathrm{m\,s^{-1}}\), so total time is less} & E1 & {States the models agree from 6 to 10, so the time saved below 6 makes the total time less.} \\
\nopagebreak & & [2] & \\
3(iv)(b) & {The same} & E1 & {States that the greatest steady speed is the same.} \\
& {\(12 > 6\), so same resistances and same maximum power at that speed} & E1 & {Explains that the steady speed of 12 exceeds 6, so the original resistances and the same maximum power apply there.} \\
\nopagebreak & & [2] & \\
\nopagebreak \SetCell[c=2]{r} Total for Question 3 & & 14 & \\
\end{msscheme}

\msQuestion{4}
\begin{msscheme}{
  hline{3-5,7-11} = {2-Z}{0.4pt, msGrid},
  hline{6,12} = {0.4pt, msGrid},
  row{12} = {font=\bfseries},
}
Question & Scheme & Marks & Notes \\
4(a) & {\(P = Fv \Rightarrow 12\,000 = 15F \Rightarrow F = 800\)} & B1 & \\
& {\(\their{800} = R + 1000 \times 9.8 \times \tfrac{1}{20}\)} & M1 A1 & {\msNote{M1}{Resolves along the slope with zero acceleration, with their driving force, \(R\) and the weight component. Condone sign slips and \(\sin\)/\(\cos\) confusion.}} \\
& {\(R = 310\)} & A1 & \\
\nopagebreak & & [4] & \\
4(b) & {\(F = \frac{12\,000}{7.5} = 1600\)} & B1 & \\
& {\(\their{1600} - \their{310} - 106 - 1600 \times 9.8 \times \tfrac{1}{20}\)\\[4pt] \(\phantom{1600} = 1600a\)} & M1 A1 & {\msNote{M1}{Forms an equation of motion for the system, or for the van alone with \(T\), with driving force, all resistances and weight components. Condone sign slips and \(\sin\)/\(\cos\) confusion.}} \\
& {\(a = 0.25\)} & A1 & \\
& {\(T - 106 - 600 \times 9.8 \times \tfrac{1}{20}\)\\[4pt] \(\phantom{T} = 600 \times \their{0.25}\)} & M1 A1 & {\msNote{M1}{Forms a second equation of motion with \(T\), for the trailer or the van, with all its forces and mass times acceleration. Condone sign slips and \(\sin\)/\(\cos\) confusion.}} \\
& {\(T = 550\)} & A1 & \\
\nopagebreak & & [7] & \\
\nopagebreak \SetCell[c=2]{r} Total for Question 4 & & 11 & \\
\end{msscheme}

\msQuestion{5}
\begin{msscheme}{
  hline{3-6} = {2-Z}{0.4pt, msGrid},
  hline{7} = {0.4pt, msGrid},
  row{7} = {font=\bfseries},
}
Question & Scheme & Marks & Notes \\
5 & {Maximum power at the greatest constant speed\\[4pt] and for the greatest acceleration:\\[4pt] \(P = Fv \Rightarrow F = \frac{64\,000}{v}\)} & E1 & {States maximum power at \(v = 20\) and at \(v = 16\), each with a reason. Both needed.} \\
& {\(\frac{64\,000}{20} = (200 + 20k) + 2000g\sin\alpha\)} & M1 A1 & {\msNote{M1}{Resolves along the slope with zero acceleration, using driving force \(\frac{64\,000}{20}\), resistance at \(v = 20\) and a weight component. Condone sign slips.}} \\
& {\(k = 101\)} & A1 & \\
& {\(\frac{64\,000}{16} - (200 + 16\their{k})\)\\[4pt] \(\phantom{\frac{64\,000}{16}} - 2000g\sin\alpha = 2000a\)} & M1 A1 & {\msNote{M1}{Applies Newton's second law along the slope, using driving force \(\frac{64\,000}{16}\), resistance \(200 + 16k\) with their \(k\), and a weight component. Condone sign slips.}} \\
& {\(a = 0.602\,\mathrm{m\,s^{-2}}\)} & A1 & \\
\nopagebreak \SetCell[c=2]{r} Total for Question 5 & & 7 & \\
\end{msscheme}

\msQuestion{6}
\begin{msscheme}{
  hline{3-6,8-9,11-12,14-15} = {2-Z}{0.4pt, msGrid},
  hline{7,10,13,16} = {0.4pt, msGrid},
  row{16} = {font=\bfseries},
}
Question & Scheme & Marks & Notes \\
6(a) & {Air resistance \(kv\): \(1200 = 600 + 15k \Rightarrow k = \ldots\)} & M1 & {Equates the driving force to \(600\) plus a resistance proportional to speed at \(v = 15\). Condone sign slips.} \\
& {\(k = 40\)} & A1 & \\
& {\(P = Fv\)\\[4pt] \(P = (600 + \their{40} \times 30) \times 30\)} & M1 & {Takes the driving force at \(30\,\mathrm{m\,s^{-1}}\) as their resistance \(600 + 30k\) and multiplies by 30.} \\
& {\(P = 1800 \times 30 = 54\,000\)} & A1~AG & {Correct \(k\) and substituted \(P = Fv\) leading without error to \(54\,000\).} \\
\nopagebreak & & [4] & \\
6(b) & {\(\frac{54\,000}{20} - (600 + \their{40} \times 20) = 1300a\)} & M1 & {Uses Newton's second law with driving force \(\frac{54\,000}{20}\), their resistance \(600 + 20k\) and mass 1300. Condone sign slips.} \\
& {\(a = 1\)} & A1 & \\
\nopagebreak & & [2] & \\
6(c) & {\(\frac{54\,000}{4} - (600 + \their{40} \times 4) = 1300a\)} & M1 & {Uses Newton's second law with driving force \(\frac{54\,000}{4}\), their resistance \(600 + 4k\) and mass 1300. Condone sign slips.} \\
& {\(a = 9.8\)} & A1 & \\
\nopagebreak & & [2] & \\
6(d) & {Acceleration at \(20\,\mathrm{m\,s^{-1}}\) reasonable;\\[4pt] at \(4\,\mathrm{m\,s^{-1}}\) unrealistically large for a van} & E1 & {Interprets both answers, with (c)'s judged too large, for example as close to \(g\). Not for ``unrealistic'' without saying which.} \\
& {Model not valid at low speeds: \(F = \frac{P}{v}\) becomes very large as \(v \to 0\)} & E1 & {Concludes the model fails at low speeds, giving the reason that \(\frac{P}{v}\) grows as \(v\) falls. Accept that full power or tyre grip cannot be sustained then.} \\
\nopagebreak & & [2] & \\
\nopagebreak \SetCell[c=2]{r} Total for Question 6 & & 10 & \\
\end{msscheme}

\msQuestion{7}
\begin{msscheme}{
  hline{3-6} = {2-Z}{0.4pt, msGrid},
  hline{7} = {0.4pt, msGrid},
  row{7} = {font=\bfseries},
}
Question & Scheme & Marks & Notes \\
7 & {\(k(14)^2 = 1000 \times 9.8 \times 0.08\ [= 784]\)} & M1 A1 & {\msNote{M1}{Equates resistance of the form \(kv^2\) at \(v = 14\) to the weight component along the slope, with no driving force.}} \\
& {\(k = 4\)} & A1 & \\
& {\(P = Fv\)\\[4pt] \(\dfrac{P}{10} = 784 + \their{4}(10)^2\)} & M1 A1 & {\msNote{M1}{Resolves up the slope with zero acceleration, using driving force \(\frac{P}{10}\) (or \(F\) with \(P = 10F\)), their \(kv^2\) at \(v = 10\) and the weight component. Condone sign slips.}\\ \msNote{A1}{Correct equation oe, such as \(F = 1184\).}} \\
& {\(P = 11\,840\)} & A1 & {awrt 11\,800 or 12\,000.} \\
& {\(R = \their{4}(12)^2 = 576\)} & B1ft & {ft their \(k\), as \(144k\).} \\
\nopagebreak \SetCell[c=2]{r} Total for Question 7 & & 7 & \\
\end{msscheme}

\msQuestion{8}
\begin{msscheme}{
  hline{3-7} = {2-Z}{0.4pt, msGrid},
  hline{8} = {0.4pt, msGrid},
  row{8} = {font=\bfseries},
}
Question & Scheme & Marks & Notes \\
8 & {\(P = Fv \Rightarrow F = \frac{9000}{12} = 750\)} & B1 & \\
& {Whole system, up the slope positive:\\[4pt] \(750 + 240 + 90 - 900g\sin\alpha = 900a\)} & M1 A1 & {\msNote{M1}{Forms an equation of motion for car and trailer with their braking force from \(\frac{P}{v}\), both resistances, the weight component and \(900a\). Condone sign slips.}\\ \msNote{A1}{Correct equation, unsimplified.}} \\
& {\(a = 0.22\)} & A1 & {Or \(-0.22\), consistent with their positive direction.} \\
& {Trailer, up the slope positive:\\[4pt] \(T + 90 - 200g\sin\alpha = 200a\)} & M1 A1 & {\msNote{M1}{Forms an equation of motion for the trailer, or the car with its \(750\,\mathrm{N}\) braking force, with the towbar force, its resistance, its weight component and mass times their \(a\). Condone sign slips.}\\ \msNote{A1}{Correct equation, in \(a\) or with \(a = 0.22\).}} \\
& {\(T = 150\)} & A1 & {Magnitude 150 N.} \\
& {Towbar in compression} & A1 & {cso, with the direction of the towbar force on the trailer or car correctly interpreted.} \\
\nopagebreak \SetCell[c=2]{r} Total for Question 8 & & 8 & \\
\end{msscheme}

\msQuestion{9}
\begin{msscheme}{
  hline{3-5,7-10} = {2-Z}{0.4pt, msGrid},
  hline{6,11} = {0.4pt, msGrid},
  row{11} = {font=\bfseries},
}
Question & Scheme & Marks & Notes \\
9(a) & {\(P = Fv \Rightarrow F = \frac{12\,000}{12} = 1000\)} & B1 & \\
& {\(\their{1000} = 1000g\sin\theta + 216 + 294\)} & M1 A1 & {\msNote{M1}{Forms an equilibrium equation along the slope for the whole system, or one for each body, with driving force, both resistances and a weight component. Condone sign slips.}\\ \msNote{A1}{Correct equation, or a correct pair of equations for the two bodies.}} \\
& {\(\sin\theta = 0.05\)} & A1 & \\
\nopagebreak & & [4] & \\
9(b) & {\(KE_0 = \tfrac12 \times 200 \times 14^2 = 19\,600\)} & B1 & \\
& {\(GPE_0 - GPE_1 = 200gs\sin\theta\ [= 98s]\)} & B1 & {GPE lost by the sledge in terms of its distance, \(200gs\sin\theta\) or \(98s\) oe.} \\
& {\(19\,600 + 200gs\sin\theta - 294s = 0 \Rightarrow s = \ldots\)} & M1 A1 & {\msNote{M1}{Forms a work--energy equation for the sledge with KE, GPE and work against resistance terms. Condone sign slips. Not for Newton's second law with constant-acceleration formulae.}\\ \msNote{A1}{Correct equation oe, with \(\sin\theta\) or 0.05.}} \\
& {\(AB = 100\,\mathrm{m}\)} & A1 & \\
\nopagebreak & & [5] & \\
\nopagebreak \SetCell[c=2]{r} Total for Question 9 & & 9 & \\
\end{msscheme}

\msQuestion{10}
\begin{msscheme}{
  hline{3-5,7-10} = {2-Z}{0.4pt, msGrid},
  hline{6,11} = {0.4pt, msGrid},
  row{11} = {font=\bfseries},
}
Question & Scheme & Marks & Notes \\
10(a) & {\(F + 5000g\sin\theta - (800 + 500) = 5000(1.2)\)\\[4pt] \(F + 5000(9.8)(0.06) - 1300 = 6000\)} & M1 A1 & {\msNote{M1}{Forms an equation of motion along the slope for the combination, or for each vehicle with a tension, with driving force, weight component and both resistances. Condone sign slips and \(\sin\)/\(\cos\) confusion.}} \\
& {\(P = Fv = \their{4360} \times 14\)} & M1 & {Uses \(P = Fv\) with their driving force and \(v = 14\), or substitutes \(F = \frac{P}{14}\) into their equation of motion.} \\
& {\(P = 61\,040\)} & A1 & {awrt \(61\,000\) oe, such as \(61.0\,\mathrm{kW}\).} \\
\nopagebreak & & [4] & \\
10(b) & {\(Pt = 61\,040 \times 7.93\)} & B1 & {Engine work \(7.93P\), with \(P\) or their \(P\) from (a).} \\
& {\(GPE_0 - GPE_1 = 5000g\,s\sin\theta = 2940s\)} & B1 & {GPE lost \(5000g\,s\sin\theta\) oe, in terms of the distance.} \\
& {\(KE_0+GPE_0+\sum W_{\text{other}}=KE_1+GPE_1\)\\[4pt] \(\tfrac12(5000)(14^2) + 2940s + \their{61\,040}(7.93) - 1300s\)\\[4pt] \(\phantom{\tfrac12(5000)(14^2)} = \tfrac12(5000)(22^2)\)} & M1 A1ft & {\msNote{M1}{Forms a work--energy equation with the KE change, engine work, GPE change and work against both resistances. Condone sign slips.}\\ \msNote{A1ft}{Correct equation with \(P\) or their \(P\) from (a).}} \\
& {\(s = \frac{720\,000 - 61\,040(7.93)}{1640}\)\\[4pt] \(AB = 144\,\mathrm{m}\)} & A1 & {awrt 144 or 140.} \\
\nopagebreak & & [5] & \\
\nopagebreak \SetCell[c=2]{r} Total for Question 10 & & 9 & \\
\end{msscheme}

\msQuestion{11}
\begin{msscheme}{
  hline{3-7} = {2-Z}{0.4pt, msGrid},
  hline{8} = {0.4pt, msGrid},
  row{8} = {font=\bfseries},
}
Question & Scheme & Marks & Notes \\
11 & {\(\sin\alpha = \frac{7}{25} \Rightarrow 2500g\sin\alpha = 6860\)} & B1 & {Correct weight component, soi by \(2500g \times \frac{7}{25}\) or 6860 in an equation.} \\
& {\(\frac{P}{12} - R - 6860 = 2500 \times 2.4\)} & M1 A1 & {\msNote{M1}{Resolves along the slope at speed 12 with \(\frac{P}{12}\), \(R\), a weight component and \(2500 \times 2.4\). Condone sign slips and \(\sin\)/\(\cos\) confusion.}} \\
& {\(\frac{P}{18} - R - 6860 = 2500 \times 0.4\)} & M1 A1 & {\msNote{M1}{Resolves along the slope at speed 18 as at speed 12, with \(\frac{P}{18}\) and \(2500 \times 0.4\).}} \\
& {\(\frac{P}{12} - \frac{P}{18} = \their{12\,860} - \their{7860} \Rightarrow P = \ldots\)} & dM1 & {Eliminates \(P\) or \(R\) from their two equations, each with \(P\) divided by the speed, and solves for one unknown.} \\
& {\(P = 180\,000\)} & A1 & {Condone an error in the weight component common to both equations, which cancels.} \\
& {\(R = 2140\)} & A1 & \\
\nopagebreak \SetCell[c=2]{r} Total for Question 11 & & 8 & \\
\end{msscheme}

\msQuestion{12}
\begin{msscheme}{
  hline{3-4,6-9,11-15} = {2-Z}{0.4pt, msGrid},
  hline{5,10,16} = {0.4pt, msGrid},
  row{16} = {font=\bfseries},
}
Question & Scheme & Marks & Notes \\
12(a) & {\(1000P = 2000 \times 20\)} & M1 & {Uses \(P = Fv\) with the driving force equal to the resistance and \(v = 20\). Condone a missing or wrong factor of 1000.} \\
& {\(P = 40\)} & A1 & {\(40\,000\) given as \(P\) scores A0.} \\
\nopagebreak & & [2] & \\
12(b) & {\(F = \frac{40\,000}{10} = 4000\)} & B1 & \\
& {\(\their{4000} - (220 + 2 \times 10^2) - 2000g\sin\alpha\)\\[4pt] \(\phantom{4000} = 2000a\)} & M1 & {Forms an equation of motion along the slope with their driving force, the refined resistance at \(v = 10\), the weight component and \(2000a\). Condone sign slips and \(\sin\)/\(\cos\) confusion.} \\
& {\(4000 - 420 - 980 = 2000a\)} & A1 & {Correct equation, oe.} \\
& {\(a = 1.3\)} & A1 & \\
\nopagebreak & & [4] & \\
12(c) & {\(F = \frac{40\,000}{U}\)} & B1 & {Driving force as their power in watts over \(U\).} \\
& {\(\their{\frac{40\,000}{U}} = 220 + 2U^2 + 2000g\sin\alpha\)} & M1 & {Equates their driving force, of the form \(\frac{k}{U}\), to \(220 + 2U^2\) plus the weight component. Condone sign slips and \(\sin\)/\(\cos\) confusion.} \\
& {\(\frac{40\,000}{U} = 1200 + 2U^2\)} & A1 & {Correct equation, oe.} \\
& {\(U^3 + 600U - 20\,000 = 0 \Rightarrow U = \ldots\)} & dM1 & {Multiplies through by \(U\) to form a cubic and solves it by any means, including a calculator.} \\
& {\(U = 20\)} & A1 & {Other roots need not be considered.} \\
\nopagebreak & & [5] & \\
\nopagebreak \SetCell[c=2]{r} Total for Question 12 & & 11 & \\
\end{msscheme}

\end{document}
